\section{Primitive actions and exact family transport}
\subsection{Primitivity without constructing a large coset action}
Let $H\leq G$ be finite. The transitive action of $G$ on its right cosets
of $H$ is primitive of degree at least two if and only if $H<G$ and
\[
 \langle H,x\rangle=G
 \quad\hbox{for every representative }x\notin H
 \hbox{ of an }(H,H)\hbox{-double coset in }G.                 \tag{*}
\]
This is the standard block--overgroup correspondence together with an
explicit double-coset test. Here is the complete argument. A block
containing the base coset $H$ is the orbit of $H$ under a subgroup
$K$ with $H\leq K\leq G$: its setwise stabilizer is $K$, and transitivity
shows the block is exactly $\{Hk:k\in K\}$. Conversely this set is a
block for each such $K$. Hence degree-at-least-two primitivity is exactly
maximal properness of $H$. If $g=h_1xh_2$, with $h_1,h_2\in H$, then
$\langle H,g\rangle=\langle H,x\rangle$. A proper intermediate subgroup
therefore contains an element outside $H$ and supplies a representative
violating (*); conversely any violating representative gives the proper
intermediate subgroup $\langle H,x\rangle$. The argument does not require
a faithful coset action or a core-free stabilizer. The case $H=G$ must be
rejected separately: its degree-one action is outside the convention and
would otherwise satisfy (*) vacuously.

The finite implementation enumerates the complete double-coset roster,
forms each indicated generated subgroup and compares its order with $|G|$.
Every negative nontrivial case records the outside element and actual
intermediate subgroup generators. For a residue it tests types in
increasing coset degree until a primitive type is found, or exhausts all
types. This preserves the existential weak-primitivity quantifier.
Applying it to all base residues is legitimate by the separately checked
flag transitivity; the residue common theorem supplies their actual
stabilizers and translates. No previously classified geometry predicate
is used by this eligibility test.
